\chapter{The three model families}
\label{ch:families}

The network class is \code{PINN} in \file{swe_pinn/physics/common.py}.
Every family accepts two scaled coordinates and returns one scalar per field, concatenated along the last axis.
One-dimensional benchmarks return \((h,u)\) or \((\eta,u)\).
The three steady square benchmarks return \((h,u,v)\).
Each field has its own tower.
There is no shared trunk.

Coordinates are scaled onto \([0,1]\) before they enter the network.
Dense-layer kernels use Glorot normal initialization and zero biases.
Hidden layers use \(\tanh\).
Output layers are linear.
Circuit angles are drawn from a normal distribution with mean zero and standard deviation \(0.1\).
\code{build_model} resets the NumPy and TensorFlow seeds to 42 before constructing the network, so a given architecture starts from the same weights in every trial and in the final training.

The legal names are \code{classical}, \code{hybrid_quantum}, and \code{full_quantum}.
Any other string raises \code{ValueError} and lists those three names.

\section{Classical}
\label{sec:classical}

Set
\begin{lstlisting}
model_family = "classical"
\end{lstlisting}
Each tower is a stack of \code{depth} hidden layers of width \code{width}, followed by a linear unit.
The circuit is not built.
\code{n_qubits}, \code{q_layers}, and \code{embedding} are ignored, and the stored embedding is \code{once}.

For one tower with two inputs and depth \(d\ge 1\), the trainable count is
\begin{equation}
\label{eq:classical}
P_{\mathrm{C}}
= 3w + (d-1)(w^{2}+w) + (w+1),
\end{equation}
where \(w\) is the width.
The first term is the map from the two coordinates into the first hidden layer.
The last term is the linear readout.
An \(m\)-field model has \(m P_{\mathrm{C}}\) parameters, all of them classical.
The banner prints \code{quantum_params=0}.

The shipped search ranges use \(d\ge 2\).
Equation~\eqref{eq:classical} is the count for that regime.
A depth of~1 is a single hidden layer: \(P_{\mathrm{C}}=4w+1\).

\section{Hybrid quantum}
\label{sec:hybrid}

Set
\begin{lstlisting}
model_family = "hybrid_quantum"
n_qubits_value = 4
q_layers_value = 2
embedding_value = "reupload"
\end{lstlisting}
The classical towers are unchanged.
After the last hidden layer of each tower the forward pass adds
\begin{equation}
\label{eq:hybrid}
z \leftarrow z +
\operatorname{decode}\!\big(
\operatorname{circuit}(\operatorname{encode}(z))
\big).
\end{equation}
The corrected state then enters the linear output layer.
Each tower has its own encoder, circuit, and decoder.
The correction is not multiplied by a trainable gate.

The encoder maps the \(w\)-dimensional hidden state to one feature per qubit.
The decoder is a \(\tanh\) dense map from the Pauli-\(Z\) expectations back to width \(w\).
With the \code{once} embedding the encoder uses \(\tanh\).
With \code{reupload} the encoder is linear, so a coordinate weight can grow past one radian.

A hybrid value of \code{n_qubits} equal to 0 is rewritten as \code{classical} before training, and the circuit is not built.
Keep \code{n_qubits_min} at 1 or more when the name is being searched, so a quantum trial actually contains a circuit.
\code{q_layers} must be at least 1 whenever the circuit is built.

For one or more qubits the count per tower is
\begin{equation}
\label{eq:hybrid-count}
P_{\mathrm{H}}
= P_{\mathrm{C}}
+ (wn_{q}+n_{q})
+ (n_{q}w+w)
+ 3n_{q}n_{L}.
\end{equation}
The three added terms are the encoder, the decoder, and the circuit rotations.
The banner splits them: encoder and decoder are classical, and the \(3n_{q}n_{L}\) angles are quantum.
An \(m\)-field model multiplies \(P_{\mathrm{H}}\) by \(m\).

\subsection{Classical warmup}

\code{hybrid_classical_warmup_fraction} is the fraction of Adam epochs, in a trial and again in the final stage, during which a hybrid tower skips the circuit.
The quantum variables still exist.
Their gradients are replaced by zeros so Adam's variable list does not change when the circuit turns on.
A fraction of 0 keeps the circuit on from epoch 1.
The fraction is clamped to the epoch budget.
The shipped \file{user_input.py} sets the fraction to 0.

Warmup applies only to \code{hybrid_quantum}.
A full-quantum model has no classical trunk to train first, and the fraction does not change that family.
After a hybrid trial the circuit is switched back on before the score lattice is evaluated.

\section{Full quantum}
\label{sec:full}

Set
\begin{lstlisting}
model_family = "full_quantum"
n_qubits_value = 4
q_layers_value = 2
embedding_value = "reupload"
\end{lstlisting}
Width and depth are stored as 0 and are not sampled.
Both \code{n_qubits} and \code{q_layers} must be at least 1.
A smaller value raises \code{ValueError} and names the field to raise.

Let \(\boldsymbol{\xi}\in\mathbb{R}^{2}\) be the two coordinates after scaling onto the unit interval.
For the tide, the dam break, and the bump, those coordinates are scaled space and time.
For the steady square problems they are scaled \(x\) and \(y\).
Each tower is
\begin{equation}
\label{eq:full}
\operatorname{readout}\!\big(
\operatorname{circuit}(\operatorname{encode}(\boldsymbol{\xi}))
\big).
\end{equation}
The encoder maps the two coordinates to \(n_{q}\) features.
The readout maps the \(n_{q}\) expectations to one scalar.
Both maps are classical dense layers.
There is no residual skip around the circuit.

The count per tower is
\begin{equation}
\label{eq:full-count}
P_{\mathrm{F}}
= 4n_{q} + 1 + 3n_{q}n_{L}.
\end{equation}
The term \(4n_{q}+1\) is the encoder (\(3n_{q}\) weights) plus the linear readout (\(n_{q}+1\)).
The banner reports those weights as classical and the \(3n_{q}n_{L}\) angles as quantum.
An \(m\)-field model has \(m P_{\mathrm{F}}\) parameters.
The family name means that the tower has no hidden perceptron.
The encoder and the readout are still classical trainable parameters.

\section{The circuit}
\label{sec:circuit}

Both quantum families use \code{QuantumLayer}.
The trainable tensor \(\theta\) has shape \code{(q_layers, n_qubits, 3)}.
That is \(3 n_{q} n_{L}\) angles per tower.

Two layouts are legal: \code{once} and \code{reupload}.
A classical model does not use either one.

\begin{description}
\item[\code{once}]
The encoder uses \(\tanh\).
\code{AngleEmbedding} places those features on the qubits with \(Y\) rotations, once, before the variational layers.
Each layer is then an input-independent \code{Rot} on every qubit, followed by the CNOT ring.

\item[\code{reupload}]
The encoder is linear.
The same features are embedded again, as \(Y\) rotations, at the start of every variational layer, and the trainable \code{Rot} follows that embedding.
\end{description}

After the rotations, a layer with more than one qubit applies CNOT gates from qubit \(i\) to \(i+1\) and a closing CNOT from the last qubit to the first.
A one-qubit layer is only the rotation.
The circuit returns the Pauli-\(Z\) expectation of every qubit.

The device is PennyLane \code{default.qubit}.
JAX vectorizes the circuit over the collocation batch with one shared set of angles, then just-in-time compiles the circuit and its derivatives.
TensorFlow still owns \(\theta\) and the tape that forms the PDE residual.
A custom gradient returns the circuit's vector--Jacobian product with respect to the features and the angles.
The feature product is differentiated again, which is the second derivative a residual needs when a slope of the network updates the angles.
The first quantum model in a process pays the compilation cost.
Later epochs of the same architecture reuse the compiled circuit.

Fix \code{embedding_value} to give every quantum tower the same layout.
Put \code{"embedding"} in \code{tune_variables} only when a study should compare layouts, and list the candidates in \code{embedding_options}.
The shipped studies fix \code{reupload} and leave \code{embedding} out of the search.

\section{Parameter counts you can check}
\label{sec:counts}

The banner printed at the start of Adam training reports classical, quantum, and total trainable parameters.
Table~\ref{tab:counts} is the count from Equations~\eqref{eq:classical}--\eqref{eq:full-count}.
A run that disagrees with this table is using a different width, depth, qubit count, layer count, or number of fields.

\begin{table}[ht]
\centering
\caption{Trainable-parameter counts for the formulas in this chapter.
Two-field models are the tide, the dam break, and the bump.
Three-field models are the steady square problems.}
\label{tab:counts}
\begin{tabular}{@{}llrrrrrr@{}}
\toprule
Family & Fields & \(w\) & \(d\) & \(n_{q}\) & \(n_{L}\) & Quantum & Total \\
\midrule
classical & 2 & 22 & 6 & 0 & 0 & 0 & 5238 \\
classical & 3 & 5 & 7 & 0 & 0 & 0 & 603 \\
classical & 3 & 1 & 2 & 0 & 0 & 0 & 21 \\
hybrid & 2 & 22 & 6 & 4 & 2 & 48 & 5690 \\
hybrid & 3 & 1 & 2 & 1 & 1 & 9 & 42 \\
full quantum & 2 & --- & --- & 4 & 2 & 48 & 82 \\
full quantum & 3 & --- & --- & 1 & 1 & 9 & 24 \\
\bottomrule
\end{tabular}
\end{table}

The three-field hybrid row is the lower edge of the parameter window in the shipped eddy study: width~1, depth~2, one qubit, and one layer.
The three-field classical row with width~5 and depth~7 is the upper edge of that same window.
Section~\ref{sec:window} explains why the window is set from those two edges.

\section{What each family is for}

Use the classical family when you want a baseline perceptron, when PennyLane is not installed, or when a search should vary only width and depth.

Use the hybrid family when the perceptron should carry the smooth part of the field and the circuit should supply a residual correction.
Warmup is available only for this family.
Width and depth still matter, and they are sampled together with the circuit size when those names are listed.

Use the full-quantum family when the field itself should be a readout of the circuit.
Width and depth do not change the model.
The capacity is the qubit count, the layer count, and the embedding.
The encoder and readout remain classical, so the banner's classical count is not zero.

All three families share the loss, the collocation grids, the Adam loop, the optional L-BFGS-B stage, and the archive format.
A comparison across families is a comparison of those networks on that shared trainer.
Chapter~\ref{ch:search} is the procedure for putting them in one study.
